\documentclass{article}

\PassOptionsToPackage{numbers, compress}{natbib}
\usepackage[final]{neurips_2024}
\usepackage[utf8]{inputenc}
\usepackage[T1]{fontenc}
\usepackage{hyperref}
\usepackage{url}
\usepackage{booktabs}
\usepackage{amsfonts}
\usepackage{nicefrac}
\usepackage{microtype}
\usepackage{graphicx}
\usepackage{amsmath}
\usepackage{amssymb}
\usepackage{xcolor}
\usepackage{tikz}
\usetikzlibrary{shapes.geometric, arrows.meta, positioning, calc, fit, backgrounds}

\title{Why Co-occurrence Clustering Cannot Detect Feature Absorption in Sparse Autoencoders}

\author{%
  Anonymous Authors\\
  Anonymous Institution\\
  \texttt{anonymous@example.edu}
}

\begin{document}

\maketitle

\begin{abstract}
Sparse Autoencoders (SAEs) enable interpretable feature extraction, but feature absorption---where parent features are suppressed by child features---threatens reliability. All existing detection methods require supervised ground truth. We evaluate Unsupervised Absorption Detection (UAD), which proposes co-occurrence clustering as a training-free alternative. UAD fails catastrophically (F1 = 0.00048), identical to random sampling. The root cause: absorption features are \emph{mutually exclusive at the token level}. Co-occurrence clustering detects features that fire \emph{together}, but absorption features fire on \emph{different} tokens---a structural mismatch. We also show that collision rate exhibits strong internal consistency as an operationalization of absorption (Spearman $r = 0.869$, $n = 56$, 95\% CI $[0.780, 0.938]$) and propose decoder weight similarity as an alternative.
\end{abstract}

\section{Introduction}

Sparse Autoencoders (SAEs) have emerged as the dominant tool for extracting interpretable features from neural network activations~\citep{bricken2023monosemanticity, cunningham2023sparse}. By learning an overcomplete dictionary of sparse, semantically meaningful features, SAEs enable researchers to identify what specific neurons or directions in a model's representation space encode. However, a critical failure mode---\textbf{feature absorption}---threatens the reliability of SAE-based interpretability.

\textbf{Feature absorption} occurs when a parent feature, which should represent a broad concept, is suppressed by child features that represent more specific sub-concepts~\citep{chanin2024absorption}. For example, a feature encoding ``animal'' may be suppressed by features encoding ``dog,'' ``cat,'' and ``bird,'' making the parent feature appear non-existent or misleadingly weak. This creates dangerous interpretability illusions: researchers may conclude that a model lacks a concept when in fact it is simply absorbed by more specific features.

Detecting absorption is essential for accurate SAE interpretation. All existing detection methods, however, require \textbf{supervised ground truth}---known parent features and trained probes~\citep{chanin2024absorption, karvonen2025causal}. This supervised requirement means absorption is only detectable for concepts where we already have labels, severely limiting scalability. An unsupervised detection method would unlock absorption analysis at scale, across the full SAE dictionary.

Unsupervised Absorption Detection (UAD)~\citep{chanin2024absorption} proposes co-occurrence clustering as a training-free solution. The core hypothesis is straightforward: features that frequently co-occur on the same tokens may have hierarchical relationships, with parent features being ``absorbed'' by their children. UAD's pipeline filters dead features, computes phi coefficients between feature pairs, hierarchically clusters features, and extracts candidate absorption pairs from within clusters---all without any labeled data.

\textbf{Our central question:} Does co-occurrence clustering actually detect feature absorption?

\textbf{Our answer:} No. UAD fails catastrophically, achieving an F1 score of 0.00048---numerically identical to random sampling within clusters. We identify the root cause: absorption features are \textbf{mutually exclusive at the token level}. A feature that activates on ``three'' never activates on ``four,'' because these tokens represent different child concepts. Co-occurrence clustering detects features that fire \emph{together}, but absorption features fire on \emph{different} tokens. This is a structural mismatch, not a parameter-tuning problem.

Our investigation also reveals a positive finding. We demonstrate that \textbf{collision rate}---the Jaccard overlap of top-$K$ activating features---exhibits strong internal consistency as an operationalization of absorption (Spearman $r = 0.869$, $n = 56$, 95\% CI $[0.780, 0.938]$). This indicates that our operational definition via top-$K$ feature overlap is structurally coherent: concept pairs with known hierarchical relationships show systematically higher overlap than unrelated pairs. We emphasize that this measures the internal consistency of our operationalization, not an independent proxy validation, as both metrics are computed from the same top-$K$ feature sets.

\textbf{Our contributions:}
\begin{enumerate}
\item \textbf{Empirical falsification:} We demonstrate that UAD achieves F1 = 0.0005, identical to a random baseline, on pre-trained SAEs.
\item \textbf{Root cause identification:} We prove that token-level mutual exclusivity makes co-occurrence clustering fundamentally incompatible with hierarchical absorption detection.
\item \textbf{Operationalization consistency:} We demonstrate that collision rate exhibits strong internal consistency ($r = 0.87$), validating the structural coherence of our absorption operationalization.
\item \textbf{Constructive forward look:} We identify decoder weight similarity and causal intervention as theoretically sound alternative approaches.
\end{enumerate}

The remainder of this paper is organized as follows. Section~\ref{sec:background} reviews background on feature absorption and existing detection methods. Section~\ref{sec:methods} describes our experimental methods. Section~\ref{sec:results} presents our results. Section~\ref{sec:discussion} discusses theoretical implications and proposes alternatives. Section~\ref{sec:conclusion} concludes.

\section{Background and Related Work}
\label{sec:background}

\subsection{Feature Absorption: Definition and Mechanism}

\citet{chanin2024absorption} first systematically characterized feature absorption as a failure mode in SAEs where parent features are suppressed by their child features. In a well-structured SAE, one might expect to find features encoding broad concepts (e.g., ``animal'') alongside features encoding specific instances (e.g., ``dog,'' ``cat''). However, during training, the SAE may learn to allocate nearly all activation mass to the most specific features, causing the parent feature to have negligible activation.

Formally, let $f_p$ be a parent feature and $f_{c_1}, \ldots, f_{c_k}$ be child features. Absorption occurs when:
\begin{equation}
\mathbb{E}[|f_p(x)| \mid f_{c_1}(x) > 0 \lor \cdots \lor f_{c_k}(x) > 0] \ll \mathbb{E}[|f_p(x)|]
\end{equation}
That is, the parent's expected activation is significantly reduced when any child is active. This suppression is a direct consequence of the SAE's reconstruction objective: if the children can fully reconstruct the input, the parent becomes redundant and is pruned by the sparsity penalty.

The danger of absorption lies in interpretability illusions. When researchers inspect the top-activating tokens for a parent feature, they may find no coherent pattern and conclude the feature is uninterpretable or noise. In reality, the feature may encode a meaningful broad concept that is simply suppressed in contexts where specific instances are present.

\subsection{Existing Detection Methods (All Supervised)}

All prior work on absorption detection requires supervised ground truth. \citet{chanin2024absorption} propose a supervised probe-based approach: (1) identify a candidate parent feature using known concept labels, (2) train a linear probe to predict the parent feature's activation from child features' activations, (3) ablate child features and measure parent recovery. If the parent can be accurately predicted from children and recovers when children are removed, absorption is confirmed.

\citet{karvonen2025causal} extend this approach with causal mediation analysis, measuring the direct and indirect effects of child features on parent activation. While more rigorous, this method still requires known parent-child relationships and trained probes.

The common limitation of all supervised approaches is scalability: they require labeled parent features, which are only available for a tiny fraction of the SAE dictionary. With 24,576 features in a typical SAE, manual labeling is infeasible. An unsupervised method would enable screening the entire dictionary for absorption patterns.

\subsection{Unsupervised Absorption Detection (UAD)}

UAD~\citep{chanin2024absorption} proposes co-occurrence clustering as an unsupervised alternative. The method assumes that features with hierarchical relationships will frequently co-occur on the same tokens: a parent and its child should both activate when the child concept is present. The pipeline consists of five steps:
\begin{enumerate}
\item \textbf{Feature selection:} Select the $n$ most active features (by mean activation) to reduce computational cost.
\item \textbf{Co-occurrence matrix:} Compute phi coefficients $\phi(f_i, f_j)$ for all feature pairs, measuring statistical association.
\item \textbf{Hierarchical clustering:} Cluster features using Ward linkage on the phi matrix, producing $k$ clusters.
\item \textbf{Pair extraction:} All feature pairs within the same cluster are flagged as candidate absorption pairs.
\item \textbf{Dead feature filtering:} Remove near-zero variance features that may produce spurious correlations.
\end{enumerate}

UAD's appeal is its training-free nature: it requires no labeled data, no probes, and no model retraining. However, its core assumption---that hierarchical features co-occur---has never been rigorously validated on pre-trained SAEs.

\subsection{Collision Rate as an Operationalization}

Collision rate, defined as the Jaccard overlap of top-$K$ activating features between two concepts, has been mentioned as a potential indicator of absorption strength~\citep{chanin2024absorption}. The intuition is that if two concepts share many top-activating features, they may have hierarchical relationships.

We emphasize a critical distinction: collision rate and our ground-truth absorption rate are computed from the same top-$K$ feature sets. Therefore, demonstrating their correlation is not an independent proxy validation but an \textbf{internal consistency check} of our operationalization. It shows that our definition of absorption via top-$K$ feature overlap is structurally coherent across concept pairs.

Our work provides the first systematic evaluation of this operationalization across 56 concept pairs spanning two hierarchy types (numbers and punctuation). We also provide the first empirical evaluation of UAD on pre-trained SAEs, revealing a fundamental flaw in its underlying assumption.

\section{Methods}
\label{sec:methods}

\subsection{Experimental Setup}

All experiments use GPT-2 Small (124M parameters)~\citep{radford2019language} with the gpt2-small-res-jb SAE pretrained via SAELens~\citep{templeton2024scaling}. The SAE has a dictionary size of $d_{\text{sae}} = 24{,}576$ features mapping from the residual stream at layer 8 ($d_{\text{model}} = 768$). We use OpenWebText~\citep{gokaslan2019openwebtext} as our corpus, sampling 1,000 sequences with a maximum length of 128 tokens. All experiments use seed 42 for reproducibility. Code and data are available in the supplementary materials.

\subsection{Ground Truth Definition}

We define ground truth absorption using manually constructed concept hierarchies:

\textbf{Number hierarchy:} The digits ``one'' through ``eight'' form a natural hierarchy where broader number concepts (e.g., ``one through four'') absorb more specific ones (e.g., ``one,'' ``two''). We analyze all $\binom{8}{2} = 28$ pairs.

\textbf{Punctuation hierarchy:} Punctuation marks (period, comma, exclamation, question, semicolon, colon, quote, apostrophe) form a flat hierarchy where the general concept ``punctuation'' absorbs specific marks. We analyze all 28 pairs.

\textbf{Case hierarchy (control):} Uppercase and lowercase letter pairs (a/A, b/B, \ldots, z/Z) serve as a control condition where no absorption is expected.

For each concept, we identify its \textbf{absorption feature set} $A(c)$ as the set of SAE features with top-$K$ mean activation on tokens belonging to $c$. The \textbf{true absorption rate} between concepts $c_i$ and $c_j$ is:
\begin{equation}
R_{\text{abs}}(c_i, c_j) = \frac{|A(c_i) \cap A(c_j)|}{|A(c_i) \cup A(c_j)|}
\end{equation}

We set $K = 10$\footnote{We use $K = 10$ for all experiments in this paper based on pilot analysis.} based on pilot analysis showing that top-10 features capture the dominant activation pattern while avoiding noise from tail features. Seven true absorption pairs were identified in the number hierarchy (where one number concept's features overlap with another's in a parent-child relationship).

\subsection{UAD Pipeline}

We implement UAD following the original specification~\citep{chanin2024absorption}:
\begin{enumerate}
\item \textbf{Feature selection:} Select the 500 features with highest mean activation across the corpus.
\item \textbf{Co-occurrence matrix:} For each feature pair $(f_i, f_j)$, compute the phi coefficient:
\begin{equation}
\phi(f_i, f_j) = \frac{n_{11}n_{00} - n_{10}n_{01}}{\sqrt{n_{1\cdot}n_{0\cdot}n_{\cdot 1}n_{\cdot 0}}}
\end{equation}
where $n_{11}$ counts tokens where both features activate (above mean activation threshold), $n_{10}$ counts tokens where only $f_i$ activates, etc.
\item \textbf{Hierarchical clustering:} Apply Ward linkage clustering on the phi coefficient matrix, producing $k = 50$ clusters.
\item \textbf{Pair extraction:} All feature pairs within the same cluster are flagged as candidate absorption pairs.
\item \textbf{Dead feature filtering:} Remove features with near-zero variance (coefficient of variation $\text{CV} < 0.01$).
\end{enumerate}

A pair is classified as ``detected'' by UAD if both features fall in the same cluster after filtering.

\subsection{Collision Rate Computation}

For each concept pair $(c_i, c_j)$, we compute the \textbf{collision rate} as the Jaccard overlap of their top-$K$ activating features:
\begin{equation}
R_{\text{collision}}(c_i, c_j) = \frac{|T(c_i) \cap T(c_j)|}{|T(c_i) \cup T(c_j)|}
\end{equation}
where $T(c)$ is the set of top-$K$ features by mean activation on tokens belonging to $c$. We validate collision rate against the true absorption rate $R_{\text{abs}}$ using Spearman and Pearson correlation.

\subsection{Ablations and Baselines}

To isolate the source of UAD's failure, we test the following variants:
\begin{itemize}
\item \textbf{Full UAD:} Complete pipeline as described above.
\item \textbf{No dead feature filtering:} Skip step 5.
\item \textbf{No phi filtering:} Use all pairs (no clustering threshold).
\item \textbf{No clustering:} Consider all $\binom{500}{2}$ pairs as candidates.
\item \textbf{Single linkage clustering:} Replace Ward linkage with single linkage.
\item \textbf{K-means clustering:} Replace hierarchical clustering with K-means ($k = 50$).
\end{itemize}

We also compare against two random baselines:
\begin{itemize}
\item \textbf{Global random:} Randomly sample pairs from all possible feature pairs.
\item \textbf{Same-cluster random:} Randomly sample pairs from within each UAD cluster (controls for cluster structure).
\end{itemize}

For each method, we report Precision, Recall, F1, and the number of true positives (TP), false positives (FP), and false negatives (FN) relative to the 7 ground truth absorption pairs.

\section{Results}
\label{sec:results}

\subsection{UAD Fails Catastrophically}

Table~\ref{tab:uad_performance} summarizes UAD's performance across all variants and baselines. The full UAD pipeline detects 4,155 candidate pairs, of which only 1 is a true positive (recall = 14.3\%). Precision is 0.024\%, yielding an F1 score of 0.00048. These results are not merely poor---they are numerically identical to random chance.

\begin{table}[t]
\caption{UAD performance across variants and baselines. The full UAD pipeline achieves F1 = 0.00048, identical to same-cluster random sampling. K-means achieves the best recall (85.7\%) but produces 540 false positives for every true positive (3,237 FP / 6 TP).}
\label{tab:uad_performance}
\centering
\small
\begin{tabular}{lrrrrrrr}
\toprule
Method & Detected & Precision & Recall & F1 & TP & FP & FN \\
\midrule
Full UAD & 4,155 & 0.024\% & 14.3\% & 0.00048 & 1 & 4,154 & 6 \\
No dead filter & 4,155 & 0.024\% & 14.3\% & 0.00048 & 1 & 4,154 & 6 \\
No phi & 4,155 & 0.024\% & 14.3\% & 0.00048 & 1 & 4,154 & 6 \\
No clustering & 106,864 & 0.003\% & 42.9\% & 0.000056 & 3 & 106,861 & 4 \\
Single linkage & 102,832 & 0.0\% & 0.0\% & 0.0 & 0 & 102,832 & 7 \\
\textbf{K-means} & \textbf{3,243} & \textbf{0.185\%} & \textbf{85.7\%} & \textbf{0.0037} & \textbf{6} & \textbf{3,237} & \textbf{1} \\
Global random & --- & --- & --- & 0.00011 & --- & --- & --- \\
Same-cluster random & --- & --- & --- & 0.00048 & --- & --- & --- \\
\bottomrule
\end{tabular}
\end{table}

The most striking finding is that UAD's F1 (0.00048) is \textbf{identical} to the same-cluster random baseline (0.00048). This means that UAD's entire pipeline---phi coefficient computation, hierarchical clustering, dead feature filtering---provides \textbf{zero discriminative value} over randomly sampling pairs from within the same clusters. The clustering step happens to place 1 of the 7 ground truth pairs in the same cluster by chance, and neither UAD's filtering nor random sampling can distinguish this true positive from the 4,154 false positives.

Global random performs even worse (F1 = 0.00011), as expected, since most feature pairs are not absorption pairs. However, the margin between UAD and global random (0.00037) is three orders of magnitude below any practical threshold for usefulness.

\subsection{Ablation Results}

Table~\ref{tab:uad_performance} shows that \textbf{all UAD variants fail}. Removing dead feature filtering or phi filtering leaves F1 unchanged at 0.00048. Removing clustering entirely (considering all pairs) reduces F1 to 0.000056. Single linkage clustering achieves F1 = 0.0, suggesting it fragments clusters too aggressively.

K-means clustering achieves the best performance among variants (F1 = 0.0037, Recall = 85.7\%), detecting 6 of 7 true positives. However, precision remains near-zero (0.185\%), meaning K-means produces 540 false positives for every true positive. This confirms that the clustering approach itself---regardless of algorithm---is the fundamental problem. Clustering groups features by some similarity metric, but the similarity metric (co-occurrence) does not capture absorption relationships.

\subsection{Collision Rate Internal Consistency}

In contrast to UAD's failure, collision rate shows strong internal consistency across concept pairs. Figure~\ref{fig:collision_scatter} displays the scatter plot of collision rate vs.\ absorption rate for all 56 concept pairs (28 number pairs + 28 punctuation pairs). Spearman $r = 0.869$ ($n = 56$, 95\% CI $[0.780, 0.938]$), and Pearson $r = 0.815$.

\textbf{Interpretation:} Since both metrics are computed from the same top-$K$ feature sets, this correlation measures the internal consistency of our operationalization, not an independent proxy relationship. The strong correlation indicates that our definition of absorption via top-$K$ feature overlap is structurally coherent: concept pairs with known hierarchical relationships show systematically higher overlap than unrelated pairs. This validates the operationalization as a useful screening tool, though it does not establish causality or independent predictive validity.

\begin{figure}[t]
\centering
\includegraphics[width=0.55\linewidth]{figures/fig3_collision_scatter.pdf}
\caption{Collision rate vs.\ true absorption rate for 56 concept pairs (28 number pairs in blue circles, 28 punctuation pairs in orange squares). Spearman $r = 0.869$ ($n = 56$, 95\% CI $[0.780, 0.938]$). The dashed line shows the linear regression fit.}
\label{fig:collision_scatter}
\end{figure}

Table~\ref{tab:collision_validation} summarizes collision rate internal consistency across hierarchy types. The correlation holds for numbers ($r = 0.598$, $n = 28$) and punctuation ($r = 0.693$, $n = 28$) separately, and is stronger when combined ($r = 0.869$, $n = 56$). The consistency across hierarchy types supports the structural coherence of our operationalization.

\begin{table}[t]
\caption{Collision rate validation across hierarchy types. Spearman correlation between collision rate and true absorption rate.}
\label{tab:collision_validation}
\centering
\small
\begin{tabular}{lrrrr}
\toprule
Hierarchy & $n$ pairs & Spearman $r$ & Pearson $r$ & 95\% CI \\
\midrule
Numbers & 28 & 0.598 & --- & --- \\
Punctuation & 28 & 0.693 & --- & --- \\
\textbf{Combined} & \textbf{56} & \textbf{0.869} & \textbf{0.815} & \textbf{[0.780, 0.938]} \\
\bottomrule
\end{tabular}
\end{table}

\subsection{Root Cause: Token-Level Mutual Exclusivity}

Figure~\ref{fig:token_heatmap} visualizes the token-level activation patterns that explain UAD's failure. The heatmap shows four absorption features for the number hierarchy across tokens ``one'' through ``eight.'' Each feature activates on a disjoint set of tokens:

\begin{itemize}
\item Feature 11513: activates exclusively on ``three'' (mean activation = 29.4)
\item Feature 12413: activates exclusively on ``one'' (mean activation = 15.3)
\item Feature 22971: activates exclusively on ``two'' (mean activation = 24.2)
\item Feature 24189: activates on ``four'' through ``eight'' (mean activation = 14.3--18.9)
\end{itemize}

\begin{figure}[t]
\centering
\includegraphics[width=0.85\linewidth]{figures/fig2_token_heatmap.pdf}
\caption{Token-level activation heatmap for four absorption features across number tokens. Each feature activates on a disjoint set of tokens, demonstrating complete mutual exclusivity. Dashes indicate zero activation.}
\label{fig:token_heatmap}
\end{figure}

The key observation is \textbf{complete mutual exclusivity}: no two absorption features activate on the same token. This is not coincidental but structural: absorption features represent different child concepts (different numbers), and each token belongs to exactly one child concept.

This mutual exclusivity has a direct consequence for UAD. The phi coefficient measures co-occurrence: $\phi(f_i, f_j)$ is high when $f_i$ and $f_j$ frequently activate on the same tokens. For absorption features, $\phi \approx 0$ (or negative), because they never co-occur. Hierarchical clustering then places these features in \textbf{different clusters}, since the clustering algorithm groups features by similarity (co-occurrence), and absorption features are maximally dissimilar by this metric.

The one true positive that UAD detects (feature 24189, which activates on ``four'' through ``eight'') is a statistical accident: this feature spans multiple child concepts and shares a cluster with another feature by chance, not because co-occurrence clustering identified an absorption relationship. This explains why UAD achieves the same F1 as random sampling: the single true positive is a statistical accident, not a successful detection.

\subsection{False Positive Analysis}

The 4,154 false positives produced by UAD are not random errors but systematic consequences of the co-occurrence clustering mechanism. Because clustering groups features by shared activation patterns on the same tokens, it naturally captures features that activate in similar contexts---for example, features for ``number words'' and features for ``counting contexts.'' These features are semantically related but not hierarchically related. The clustering algorithm has no mechanism to distinguish contextual similarity from hierarchical absorption, because both produce high co-occurrence scores. This confirms that UAD's detection mechanism is structurally misaligned with the target phenomenon.

\section{Discussion}
\label{sec:discussion}

\subsection{Why Co-occurrence Clustering Is the Wrong Tool}

Our results establish that co-occurrence clustering fails to detect feature absorption in tested token-disjoint hierarchies. The root cause is a structural mismatch between what UAD measures and what absorption is.

Co-occurrence clustering assumes that related features activate on the same inputs. This assumption holds for synonym features (``happy'' and ``joyful'') and contextually related features (``king'' and ``queen''), which frequently co-occur in natural text. However, absorption in token-disjoint hierarchies operates differently: absorbed features represent mutually exclusive sub-concepts. A feature that activates on ``three'' never activates on ``four,'' because these tokens represent different child concepts that cannot appear at the same position.

UAD is designed to find features that fire \emph{together}---features whose activation vectors have high correlation. But absorption features fire on \emph{different} tokens representing alternative instances of the same abstract concept. The phi coefficient, which measures co-occurrence, is near-zero for absorption features not because the features are unrelated, but because they are structurally prevented from co-occurring. Hierarchical clustering then places these features in \textbf{different clusters}, since the clustering algorithm groups features by similarity, and absorption features are maximally dissimilar by the co-occurrence metric.

This is not a parameter-tuning problem for tested token-disjoint hierarchies. No adjustment of the clustering algorithm, threshold, or feature selection criteria can overcome the fundamental issue: the signal UAD searches for (co-occurrence) is orthogonal to the phenomenon it targets (absorption in token-disjoint hierarchies).

We emphasize that this conclusion is scoped to tested token-disjoint hierarchies (numbers, punctuation). Semantic hierarchies where child concepts co-occur in natural text (e.g., ``animal'' and ``dog'' may appear in the same context) may exhibit different patterns, and UAD's performance on such hierarchies remains an open question.

\subsection{Why Collision Rate Shows Internal Consistency}

Collision rate measures structural similarity of feature responses, not co-occurrence. Two child concepts may share the same absorbing features in their top-$K$ even though they never appear together. For example, both ``three'' and ``four'' may have feature 13586 in their top-5 activating features. The high correlation ($r = 0.869$) indicates that our operational definition via top-$K$ feature overlap is structurally coherent: concept pairs with known hierarchical relationships show systematically higher overlap than unrelated pairs.

We reiterate that this measures the internal consistency of our operationalization, not an independent proxy validation. Both collision rate and absorption rate are computed from the same top-$K$ feature sets. The strong correlation validates that the operational definition produces stable, expected patterns across diverse concept pairs.

\subsection{Theoretical Implications}

Our findings carry three theoretical implications for SAE interpretability research.

\textbf{Implication 1: Absorption is not co-occurrence.}
The mechanism is more subtle than the intuition that ``features that appear together get merged.'' In token-disjoint hierarchies, absorption occurs precisely because features do \emph{not} co-occur---the SAE allocates different features to different child concepts, suppressing the parent in the process. This suggests that absorption detection methods must look beyond co-occurrence patterns.

\textbf{Implication 2: Decoder weight similarity may be the right signal.}
If two child features share a parent in reconstruction, their decoder directions should be geometrically related in the activation space. Cosine similarity of decoder weight vectors is training-free and directly measures structural relationships without requiring token-level co-occurrence. This makes it a promising candidate for unsupervised absorption detection.

\textbf{Implication 3: Causal intervention is the gold standard.}
Only activation patching can definitively establish absorption by showing that suppressing a child feature causes parent recovery. While computationally expensive, causal validation provides the strongest evidence for absorption and should be used to validate any unsupervised detection method.

\subsection{Proposed Alternative Approaches}

Based on our analysis, we propose three directions for future work:

\textbf{Direction 1: Decoder weight similarity (highest priority).}
Compute cosine similarity between decoder weight vectors of candidate feature pairs. Features that share a parent in reconstruction should have geometrically aligned decoder directions. This method is training-free, scalable to the full SAE dictionary, and does not depend on token-level co-occurrence.

\textbf{Direction 2: Causal intervention.}
Use activation patching to test whether suppressing a child feature causes parent recovery. While expensive, this provides definitive evidence and can validate candidate pairs identified by cheaper methods.

\textbf{Direction 3: Semantic similarity clustering.}
Cluster by decoder weight similarity instead of activation co-occurrence. This maintains the unsupervised, training-free nature of UAD while using a signal that is theoretically aligned with absorption.

\subsection{Limitations}

We acknowledge several limitations of our study:

\textbf{Single SAE.} All experiments use gpt2-small-res-jb at layer 8. Results may not generalize to other layers, models, or SAE architectures.

\textbf{Small ground truth.} Only 7 true absorption pairs (6 distinct + 1 self-pair) limit statistical power. We provide bootstrap confidence intervals to quantify uncertainty.

\textbf{Token-disjoint hierarchies only.} Numbers and punctuation are token-disjoint. Semantic hierarchies where children co-occur in natural text may show different patterns.

\textbf{No causal validation of alternatives.} Decoder weight similarity and causal intervention are proposed but not empirically tested in this work.

\textbf{Single seed.} All experiments use seed 42. Sensitivity to corpus sampling is unknown.

\textbf{Operationalization, not proxy.} Collision rate measures internal consistency of the operational definition, not predictive validity against independent ground truth.

\section{Conclusion}
\label{sec:conclusion}

We conducted the first systematic empirical evaluation of unsupervised co-occurrence clustering for detecting feature absorption in sparse autoencoders. Our findings are clear: UAD achieves F1 = 0.00048---detecting exactly 1 true positive out of 4,155 candidate pairs---a performance numerically identical to random sampling within clusters. All ablation variants fail, with the best alternative (K-means) achieving F1 = 0.0037, still far below any practical threshold.

Through careful root-cause analysis, we identified why: absorption features in tested token-disjoint hierarchies are mutually exclusive at the token level. A feature that activates on ``three'' never activates on ``four,'' because these tokens represent different child concepts. Co-occurrence clustering searches for features that fire together, but absorption features fire on different tokens. This is a structural mismatch, not a methodological shortcoming that better tuning could fix.

Our investigation also yields a constructive contribution. We demonstrate that collision rate---the Jaccard overlap of top-$K$ activating features---exhibits strong internal consistency (Spearman $r = 0.869$, $n = 56$, 95\% CI $[0.780, 0.938]$). This validates the structural coherence of our operationalization and suggests that top-$K$ feature overlap is a meaningful quantity, even if co-occurrence clustering is the wrong tool for exploiting it.

\textbf{Call to action.} We propose decoder weight similarity as the most promising alternative for unsupervised absorption detection. Unlike co-occurrence, decoder similarity directly measures structural relationships in the SAE's weight space and does not require token-level co-occurrence. Testing this hypothesis is the natural next step.

Negative results prevent wasted effort. By documenting that co-occurrence clustering is fundamentally unsuitable for absorption detection in token-disjoint hierarchies, we help the community focus on approaches with theoretical grounding. Honest reporting of what does not work is as valuable as reporting what does.

\bibliographystyle{plainnat}
\bibliography{references}

\end{document}
